% 核心结论：电阻电容和电感的阻抗与导纳随角频率变化的趋势
% 图原型：纯几何／对比。物理网络按本题节点、端口和符号构建。
\documentclass[tikz,border=7pt]{standalone}
\usepackage{fontspec,xeCJK,amsmath,amssymb}
\IfFontExistsTF{Times New Roman}{\setmainfont{Times New Roman}}{\setmainfont{TeX Gyre Termes}}
\IfFontExistsTF{Heiti SC}{\setCJKmainfont[AutoFakeBold=1.5]{Heiti SC}}{\setCJKmainfont[AutoFakeBold=1.5]{FandolHei-Regular}}
\usetikzlibrary{arrows.meta,calc,angles,quotes}
\definecolor{ink}{HTML}{243745}
\definecolor{blue}{HTML}{275DAB}
\definecolor{green}{HTML}{177C63}
\definecolor{red}{HTML}{BD4935}
\definecolor{purple}{HTML}{715393}
\begin{document}
\begin{tikzpicture}[font=\fontsize{9.5}{12}\selectfont,text=ink,
 v/.style={line width=1.1pt,-{Stealth[length=2.4mm,width=1.5mm]}},
 axis/.style={ink!65,line width=.45pt,-{Stealth[length=1.7mm,width=1.1mm]}},
 aux/.style={ink!40,densely dashed,line width=.5pt},
 lab/.style={fill=white,inner sep=2pt},
 card/.style={draw=blue!25,fill=blue!3,rounded corners=3pt,text width=6cm,align=left,inner sep=10pt}]
\path[use as bounding box] (-8,-3.35) rectangle (8,3.65);
\node[font=\fontsize{10.5}{13}\selectfont\bfseries,text width=15cm,align=center] at (0.00000,3.22000) {M6-H04　电阻电容和电感的阻抗与导纳随角频率变化的趋势};
\begin{scope}[shift={(-4,0)}]
\draw[axis] (-2.50000,-1.80000) -- (2.72000,-1.80000);
\draw[axis] (-2.50000,-1.85000) -- (-2.50000,1.97000);
\node[] at (2.60000,-2.55000) {$\omega\;(rad/s)$};
\node[] at (-2.25000,2.12000) {$|X|/\Omega$};
\draw[blue,thick] plot[domain=100:1000,samples=250] ({-2.5+0.005666666666666666*(\x-100)},{-1.8+0.03318181818181819*((.1*\x)-(0))});
\draw[green,thick] plot[domain=100:1000,samples=250] ({-2.5+0.005666666666666666*(\x-100)},{-1.8+0.03318181818181819*((10000/\x)-(0))});
\draw[ink] (-2.50000,-1.86000) -- (-2.50000,-1.74000);
\node[] at (-2.50000,-2.20000) {$100$};
\draw[ink] (2.60000,-1.86000) -- (2.60000,-1.74000);
\node[] at (2.60000,-2.20000) {$1000$};
\draw[ink] (-2.56000,-1.80000) -- (-2.44000,-1.80000);
\node[anchor=east] at (-2.72000,-1.80000) {$0$};
\draw[ink] (-2.56000,1.85000) -- (-2.44000,1.85000);
\node[anchor=east] at (-2.72000,1.85000) {$110$};
\end{scope}
\begin{scope}[shift={(4,0)}]
\node[card] at (0.00000,1.15000) {$|Z_L|\uparrow,\quad|Z_C|\downarrow$};
\end{scope}
\begin{scope}[shift={(4,0)}]
\node[card] at (0.00000,-1.15000) {$|Y_L|\downarrow,\quad|Y_C|\uparrow$};
\end{scope}
\node[text width=15cm,align=center] at (0.00000,-3.00000) {复平面与实际电路分区；相量统一采用有效值和 $e^{j\omega t}$。};

\end{tikzpicture}\end{document}